% chapters/06_weiterfuehrende_arbeiten.tex

\chapter{Weiterführende Arbeiten}
\label{chap:weiterfuehrend}

Vier Stränge weisen über den Projektkern hinaus: die \textbf{Lokomotion}
(recherchiert, nicht implementiert), das \textbf{Co-Training} auf echten und
gerenderten Bildern (Werkzeugkette vollständig, Lauf ausstehend), die
\textbf{generative Video-Augmentierung} mit Cosmos-Transfer (paralleler
Teamstrang, in Entwicklung) und das \textbf{RL-Fine-Tuning} (Pipeline lief
end-to-end durch, Lernwirkung unbelegt). Die ausführlichen Konzepte liegen in
\doc{docs/weiterfuehrend/} und \doc{docs/training/co-training.md}; dieses
Kapitel fasst Argumentation, Entscheidungen und Stand zusammen.

\section{Lokomotion}
\label{sec:locomotion}
\label{sec:loco-pfade}

Lokomotion ist kein Schalter im bestehenden Modell, sondern eine zweite,
eigenständige Steuerebene. Drei voneinander unabhängige Blocker verhindern sie
derzeit (\cref{tab:loco-blocker}). Die Kinematik ist allerdings vollständig
vorhanden: Das USD-Asset wird aus dem 29-\ac{dof}-\ac{urdf} mit freier Basis
gebaut; \texttt{fix\_root\_link=True} überschreibt dies nur zur Laufzeit.

\begin{table}[htb]
  \centering
  \begin{tabularx}{\linewidth}{lX}
    \toprule
    \textbf{Ebene} & \textbf{Blocker} \\
    \midrule
    Simulation & \texttt{fix\_root\_link=True} verschweißt den Torso mit der
      Welt (bewusste Vereinfachung für die Manipulations-Eval); es gibt weder
      Boden noch Bein-Aktuatoren. \\
    Modell & Der Aktionsraum umfasst nur 28~\acp{dof} (Arme + Hände); der
      Aktionskopf kann keine Beinbefehle erzeugen. \\
    Daten  & Alle \texttt{G1\_Dex3}-Datensätze wurden an einem
      \emph{stehenden} Roboter teleoperiert; selbst ein perfektes Modell
      könnte daraus keine Beinaktionen lernen. \\
    \bottomrule
  \end{tabularx}
  \caption{Die drei unabhängigen Lokomotions-Blocker.}
  \label{tab:loco-blocker}
\end{table}

NVIDIA adressiert Manipulation und Fortbewegung im selben Humanoiden über
einen \emph{entkoppelten} \acf{wbc}~\cite{groot2025}: Eine in Isaac~Lab
trainierte RL-Politik hält mit Beinen und Taille die Balance und folgt
einfachen Geschwindigkeitsbefehlen $(v_x, v_y, \omega_{\text{yaw}})$, während
der Oberkörper von \ac{ik} bzw.\ der hier feinabgestimmten
Manipulations-Politik gesteuert wird. Einsatzfähige RL-Lokomotions-Stacks für
den G1 existieren (\texttt{unitree\_rl\_gym}, \texttt{unitree\_rl\_lab},
NVIDIAs \texttt{GR00T-WholeBodyControl}); zu beachten ist ein
Versionskonflikt -- \texttt{unitree\_rl\_lab} verlangt
Isaac~Lab~2.3/Isaac~Sim~5.1, die Sim-Eval dieses Projekts läuft auf einer
anderen Version. Die Recherche
(\doc{docs/weiterfuehrend/lokomotion-recherche.md}) empfiehlt einen gestuften
Weg: zunächst eine reine Sim-Lauf-Demo mit einer Stock-Politik (validiert
Asset, Boden und Aktuatoren), dann die entkoppelte Integration --
\texttt{fix\_root\_link} lösen, Boden und Schwerkraft ergänzen, Bein- und
Taillen-Aktuatoren hinzufügen, fertige Lauf-Politik laden, Geschwindigkeits-
befehle zunächst skriptgesteuert -- und erst langfristig eine unified
Whole-Body-\ac{vla} mit latenten Aktions-Tokens, die neue
Ganzkörper-Trainingsdaten erfordert. Das reale Unitree-SDK spiegelt dieselbe
Entkopplung (\texttt{LocoClient} für Balance, \texttt{arm\_sdk} für die
GR00T-Targets).

\section{Co-Training auf echten und gerenderten Bildern}
\label{sec:cotrain}

Die Diskussion in \cref{sec:frozen-vs-tuned} endete mit einer präzisen
Bedingung: Das Mittrainieren des Vision-Encoders hilft gegen den Domain-Gap
nur, wenn der Encoder im Training tatsächlich Simulationsbilder sieht --
in den Läufen~2 und~3 war das nicht erfüllt. Der zweite Hebel, das Angleichen
der Sim-Optik (Visual Matching), wurde in \cref{sec:neukalibrierung}
ausgeschöpft; die kritische Handgelenkkamera verfehlte das Zielkriterium
dennoch knapp. Damit bleibt der dritte Weg: die Simulationsbilder \emph{in den
Trainingssatz aufnehmen}.

Der Renderer \doc{Simulation/g1\_dex3\_sim/render\_cotrain\_dataset.py} spielt
dazu die \emph{echten} Datensatz-Aktionen in der Simulation ab und zeichnet
die vier Policy-Kameras in der kalibrierten Auflösung auf; das Ergebnis ist
ein LeRobot-Datensatz aus Paaren von \emph{simuliertem Bild} und \emph{realer
Aktion}. Das Training mischt beide Quellen in einstellbarem Verhältnis
(empfohlen: ein Viertel gerenderte Samples); die zurückgehaltenen
Testepisoden werden nie gerendert, damit der Split intakt bleibt. Das
Verfahren steht und fällt mit der Korrektheit der Simulationsszene -- ein an
falscher Stelle gerenderter Würfel erzeugt ein Bild-Aktions-Paar, das dem
Modell etwas Falsches beibringt. Genau deshalb war die geometrische
Rekonstruktion aus \cref{sec:geometrie} die Voraussetzung dieses Schritts,
und aus dem Widerruf des Vorzeichen-Fixes (\cref{sec:vorzeichen-widerruf})
folgt, dass ein zuvor erzeugter gerenderter Datensatz neu zu erzeugen ist.
Die Werkzeugkette ist vollständig und dokumentiert
(\doc{docs/training/co-training.md}); der Trainingslauf auf dem gemischten
Datensatz steht zum Redaktionsschluss aus. Er ist der nächste inhaltliche
Schritt des Projekts und zugleich Voraussetzung für einen aussagekräftigen
RL-Lauf (\cref{sec:rl-stand}).

Parallel dazu verfolgt das Team einen \emph{generativen} Ansatz
(federführend Maximilian Berger, \cref{sec:eigenleistung}): NVIDIAs Video-Weltmodell
\textbf{Cosmos-Transfer} stylt die realen Trainingsvideos um und erhöht so
die visuelle Vielfalt der Trainingsverteilung -- Domain Randomization ohne
neue Teleoperation, komplementär zum gerenderten Datensatz. Die Pipeline ist
als eigener Container aufgesetzt, zum Redaktionsschluss in Entwicklung und
nicht Teil der hier bewerteten Ergebnisse.

\section{Reinforcement-Learning-Fine-Tuning}
\label{sec:rl-weiterfuehrend}

\subsection{Motivation und Verfahren}
\label{sec:rl-verfahren}

Das bisherige Training ist reines \acl{bc}: Es optimiert
\enquote{Aktion nachahmen}, nicht \enquote{Aufgabe lösen}. \acl{rl}
optimiert direkt auf den Aufgabenerfolg und exploriert Zustände jenseits der
Demonstrationen -- genau die Korrekturen, die \ac{bc} strukturell fehlen
(\cref{sec:imitation-learning}); $\pi_{\mathrm{RL}}$ hebt die mittlere
Erfolgsrate flow-basierter \acp{vla} etwa von~$41{,}6\,\%$
auf~$85{,}7\,\%$~\cite{pirl2025}. Gemeint ist das Fine-Tuning der
\ac{vla}-Politik, nicht das Low-Level-RL des Ganzkörper-Controllers aus
\cref{sec:locomotion}. Die zentrale Hürde benennt \cref{sec:flow-matching}:
Klassisches \ac{ppo} benötigt die Aktions-Likelihood, die der
Flow-Matching-Kopf nicht liefert; die Forschung 2025 bietet passende
Verfahren (\cref{tab:rl-verfahren}).

\begin{table}[htb]
  \centering
  \begin{tabularx}{\linewidth}{lX}
    \toprule
    \textbf{Verfahren} & \textbf{Kerntrick} \\
    \midrule
    $\pi_{\mathrm{RL}}$ (Flow-SDE) & \ac{ode}$\to$SDE für handhabbare
      Likelihood, dann \ac{ppo}; architektonisch am nächsten an
      GR00T~\cite{pirl2025} \\
    FPO & likelihood-freier Proxy aus dem Flow-Matching-Loss,
      \ac{ppo}-Clipping \\
    ReinFlow & lernbares Rauschnetz $\to$ Markov-Kette mit exakter
      Likelihood \\
    \bottomrule
  \end{tabularx}
  \caption{RL-Verfahren für flow-basierte Politiken
    (\doc{docs/weiterfuehrend/reinforcement-learning-plan.md}).}
  \label{tab:rl-verfahren}
\end{table}

\subsection{Umsetzung und Stand}
\label{sec:rl-stand}

Umgesetzt wurde \textbf{FPO}: Sein likelihood-freier Proxy verwendet den
vorhandenen Flow-Matching-Loss unverändert wieder und greift am wenigsten in
die GR00T-Trainings-API ein. Der Trainer (\texttt{rl\_finetune.py}) tunt
ausschließlich den Aktionskopf, regularisiert per KL-Term gegen den
\ac{bc}-Checkpoint und bezieht seinen geformten Reward aus der entsprechend
erweiterten Isaac-Lab-Umgebung, die als \texttt{DirectRLEnv} bereits
RL-tauglich angelegt war. Den Hardware-Konflikt -- bildbasiertes RL braucht
RT-Core-Rendering \emph{und} Trainings-FLOPs zugleich, während A100/H100
keine RT-Cores besitzen (\cref{sec:eval-methodik}) -- löste der
institutseigene Blackwell-Server (zwei RTX~PRO~6000, je 96~GB; das
KL-Referenzmodell läuft auf der zweiten GPU); der Weg dorthin erforderte die
Migration auf Isaac~Sim~6.0 (\cref{sec:neukalibrierung}). Am
\textbf{08.08.2026} terminierte erstmals eine vollständige Trainingsiteration
-- Rollout $\to$ \ac{gae} $\to$ FPO-Update $\to$ Checkpoint -- fehlerfrei
(Smoke-Test mit zwei parallelen Umgebungen;
\doc{docs/weiterfuehrend/rl-anleitung.md}).

Damit ist die \emph{Machbarkeit} der Pipeline belegt, nicht ihre
\emph{Lernwirkung}. Der zwischenzeitlich als Blocker geführte
Greif-Physik-Befund ist ausgeräumt (\cref{sec:grasp-befund}); an seine Stelle
treten zwei Voraussetzungen: Erstens muss die Politik überhaupt einen
hinreichenden Griff kommandieren, damit der geformte Reward einen Gradienten
liefert -- bei rund einem Viertel der Demonstrations-Spanne wird der
Stapel-Term praktisch nie erreicht; genau hier setzt das Co-Training an, das
dem RL-Lauf deshalb vorausgehen sollte. Zweitens sind die Sim-Messungen nach
den Geometriekorrekturen aus \cref{sec:geometrie} zu wiederholen, damit ein
belastbarer Nullpunkt existiert (\cref{sec:vorbehalt-greifzahlen}).
Konzeptionell löst RL zugleich den Domain-Gap auf: Trainiert wird in genau
der Simulation, in der evaluiert wird -- der Distribution-Shift verschwindet,
und der Reward erzwingt Handeln statt Einfrieren.
